%% Copernicus Publications Manuscript Preparation Template for LaTeX Submissions
%%
%% Journal: Geoscientific Model Development (GMD)
%%
%% To compile this manuscript, you need the Copernicus LaTeX class files.
%% Download the latest template package from:
%% https://publications.copernicus.org/for_authors/manuscript_preparation.html
%%
%% Place copernicus.cls and copernicus.bst in this directory before compiling.
%%
%% Compile with:
%%   pdflatex ms.tex
%%   bibtex ms
%%   pdflatex ms.tex
%%   pdflatex ms.tex

\documentclass[gmd, manuscript]{copernicus}

\begin{document}

\title{ICON Validation with Observations: a Python-based framework for
  systematic evaluation of the ICON atmospheric model}

%% Author list and affiliations
%% Use \Author[affil_number]{First name}{Last name}
%% Use \affil[number]{Institution, City, Country}

\Author[1]{}{[Author 1]}
\Author[1,2]{}{[Author 2]}

\affil[1]{[Institution 1, City, Country]}
\affil[2]{[Institution 2, City, Country]}

%% Running title (max 60 characters) and running author (max 30 characters)
\runningtitle{ICON Validation with Observations}
\runningauthor{[Author 1] et al.}

%% Correspondence details
\correspondence{[Author 1] ([email@institution.edu])}

%% Dates -- leave blank until submission/peer review
\received{}
\pubdiscuss{}
\revised{}
\accepted{}
\published{}

%% These dates will be inserted by Copernicus Publications during production
\firstpage{1}

\maketitle

%% ============================================================
%% ABSTRACT
%% ============================================================
\begin{abstract}
The ICOsahedral Non-hydrostatic (ICON) model is a next-generation
atmospheric general circulation model jointly developed by the German
Weather Service (DWD) and the Max Planck Institute for Meteorology (MPI-M).
Systematic validation of ICON simulations against observations is essential
for assessing model performance and guiding further model development.
Here we present \texttt{icon\_validation}, an open-source Python framework
that provides a standardised, reproducible workflow for evaluating ICON
model output against a variety of observational datasets.
The framework supports comparisons of standard atmospheric variables
including temperature, precipitation, wind speed and cloud cover across
multiple temporal and spatial scales.
We describe the design principles of the software, illustrate its use with
selected validation examples, and discuss its role within a broader model
development cycle.
\end{abstract}

%% ============================================================
%% INTRODUCTION
%% ============================================================
\introduction

Numerical weather and climate models require continuous evaluation against
observations to identify systematic biases, diagnose model deficiencies and
track the impact of model improvements over successive development cycles.
The ICOsahedral Non-hydrostatic (ICON) model \citep{zaengl2015icon} is
used operationally by the German Weather Service (DWD) for numerical weather
prediction and by the Max Planck Institute for Meteorology (MPI-M) for
climate research.
Given the wide range of applications and the rapid development pace of ICON,
there is a clear need for a flexible, community-accessible validation
framework.

While general model evaluation tools such as ESMValTool
\citep{righi2020esmvaltool} and ICARE-NG exist, dedicated workflows that
are tightly integrated with the ICON data format and metadata conventions
can greatly reduce the overhead for ICON developers and users.
The \texttt{icon\_validation} package aims to fill this gap by providing
a lightweight, extensible framework built on well-established scientific
Python libraries (NumPy, xarray, Matplotlib, cartopy).

The remainder of this paper is structured as follows.
Section~\ref{sec:description} describes the software design and
implementation.
Section~\ref{sec:validation} presents selected validation results
demonstrating the capabilities of the framework.
Section~\ref{sec:conclusions} summarises the main conclusions and
outlines planned future developments.

%% ============================================================
%% SOFTWARE DESCRIPTION
%% ============================================================
\section{Software description}
\label{sec:description}

\subsection{Design principles}

The \texttt{icon\_validation} framework follows several design principles
that are intended to maximise usability, reproducibility and
maintainability:

\begin{itemize}
  \item \textbf{Modularity.}  Individual validation tasks (data ingestion,
        regridding, statistical analysis, plotting) are implemented as
        independent, reusable components.
  \item \textbf{Reproducibility.}  All validation workflows are driven by
        plain-text configuration files so that results can be reproduced
        from a single command.
  \item \textbf{Open standards.}  The framework relies exclusively on
        community-standard file formats (NetCDF, GRIB2) and Python
        packages available through conda-forge.
  \item \textbf{Scalability.}  Lazy loading via xarray and Dask enables
        processing of large multi-year datasets without loading all data
        into memory simultaneously.
\end{itemize}

\subsection{Framework architecture}

\subsubsection{Input data handling}

ICON produces output on an unstructured triangular grid.
The framework reads ICON NetCDF output directly using xarray and applies
an optional remapping step to a regular latitude--longitude grid using
CDO (Climate Data Operators) or the \texttt{xesmf} package prior to
comparison with observational data.

\subsubsection{Observational reference datasets}

The following observational and reanalysis datasets are currently supported:

\begin{itemize}
  \item ERA5 reanalysis \citep{hersbach2020era5}
  \item GPCP monthly precipitation \citep{adler2018gpcp}
  \item CERES-EBAF radiation budget \citep{loeb2018ceres}
  \item IGRA radiosonde archive
  \item SYNOP surface station data
\end{itemize}

Support for additional datasets can be added by implementing the
\texttt{ObsDataset} base class interface.

\subsubsection{Statistical diagnostics}

For each variable and dataset pair the framework computes a standard set of
scalar metrics including the mean bias, root-mean-square error (RMSE),
pattern correlation coefficient and normalised standard deviation.
These metrics are assembled into a Taylor diagram \citep{taylor2001taylor}
for visual summary.
Regional and seasonal breakdowns are also available.

\subsubsection{Visualisation}

Validation results are rendered to publication-quality figures using
Matplotlib and cartopy.
A dedicated HTML report is optionally generated, aggregating all diagnostic
plots for easy browsing.

\subsection{Installation and dependencies}

The package is available on GitHub
(\url{https://github.com/ClimLabTools/icon_validation}) and can be
installed using \texttt{pip}:
\begin{verbatim}
pip install icon-validation
\end{verbatim}
or via conda:
\begin{verbatim}
conda install -c conda-forge icon-validation
\end{verbatim}

The main dependencies are Python ($\geq$\,3.9), NumPy, xarray, Matplotlib,
cartopy, and CDO (optional, for remapping).

%% ============================================================
%% VALIDATION EXAMPLES
%% ============================================================
\section{Validation examples}
\label{sec:validation}

\subsection{Experimental setup}

To illustrate the capabilities of \texttt{icon\_validation} we evaluate
a 10-year ICON-AES atmosphere-only simulation at approximately 80\,km
horizontal resolution.
The simulation was forced with prescribed sea surface temperatures from
the HadISST dataset and covers the period 2010--2019.

\subsection{Near-surface temperature}

Figure~\ref{fig:temp_bias} shows the annual-mean near-surface (2\,m)
temperature bias of ICON relative to ERA5.
The overall global mean bias is [X]\,K, within the range of previously
published ICON results \citep{giorgetta2018icon}.

\begin{figure}[t]
  \includegraphics[width=\columnwidth]{figures/temp_bias.pdf}
  \caption{Annual-mean 2\,m temperature bias (ICON minus ERA5, K) for
           the 2010--2019 period.}
  \label{fig:temp_bias}
\end{figure}

\subsection{Precipitation}

The spatial pattern of annual-mean precipitation is evaluated against the
GPCP climatology (Fig.~\ref{fig:precip_bias}).
ICON captures the major precipitation regimes, including the
inter-tropical convergence zone (ITCZ) and the storm track regions,
though a dry bias is apparent over the Maritime Continent.

\begin{figure}[t]
  \includegraphics[width=\columnwidth]{figures/precip_bias.pdf}
  \caption{Annual-mean precipitation bias (ICON minus GPCP,
           mm\,day$^{-1}$) for the 2010--2019 period.}
  \label{fig:precip_bias}
\end{figure}

\subsection{Taylor diagram summary}

Figure~\ref{fig:taylor} shows a Taylor diagram summarising the
performance of ICON for a range of atmospheric variables.
[Discuss Taylor diagram results here.]

\begin{figure}[t]
  \includegraphics[width=\columnwidth]{figures/taylor_diagram.pdf}
  \caption{Taylor diagram for key atmospheric variables evaluated
           against ERA5 and GPCP.}
  \label{fig:taylor}
\end{figure}

%% ============================================================
%% CONCLUSIONS
%% ============================================================
\conclusions

We have presented \texttt{icon\_validation}, a Python-based framework for
systematic evaluation of the ICON atmospheric model against observations.
The framework provides a modular, reproducible workflow covering data
ingestion, regridding, statistical analysis and visualisation.
Selected validation examples demonstrate that the framework correctly
identifies known ICON biases and is capable of producing publication-quality
diagnostics.

Future development priorities include:
\begin{itemize}
  \item Extension to ocean and land-surface components of the ICON model
        family (ICON-O, ICON-Land);
  \item Integration with the ICON model test suite for continuous
        performance monitoring;
  \item Support for ensemble validation and multi-model comparison;
  \item A web-based interactive dashboard.
\end{itemize}

The source code, documentation and issue tracker are hosted on GitHub at
\url{https://github.com/ClimLabTools/icon_validation}.
Community contributions are welcome via the GitHub pull-request workflow.

%% ============================================================
%% CODE AVAILABILITY
%% ============================================================
\codeavailability{The \texttt{icon\_validation} source code is freely
available under the MIT licence at
\url{https://github.com/ClimLabTools/icon_validation}
(last access: [date]).
The version used to produce the results in this paper is archived at
\url{https://doi.org/[ZENODO_DOI]} \citep{icon_validation_code}.}

%% ============================================================
%% DATA AVAILABILITY
%% ============================================================
\dataavailability{ERA5 data are available from the Copernicus Climate
Change Service Climate Data Store
(\url{https://cds.climate.copernicus.eu}, \citealt{hersbach2020era5}).
GPCP precipitation data are available from
\url{https://psl.noaa.gov/data/gridded/data.gpcp.html}
\citep{adler2018gpcp}.
CERES-EBAF data are available from
\url{https://ceres.larc.nasa.gov/data/} \citep{loeb2018ceres}.}

%% ============================================================
%% AUTHOR CONTRIBUTIONS
%% ============================================================
\authorcontribution{[Author 1] conceived the study and led the software
development.
[Author 2] contributed to the observational data handling and validation
diagnostics.
All authors contributed to writing and revising the manuscript.}

%% ============================================================
%% COMPETING INTERESTS
%% ============================================================
\competinginterests{The authors declare that they have no conflict of
interest.}

%% ============================================================
%% ACKNOWLEDGEMENTS
%% ============================================================
\begin{acknowledgements}
The authors thank [names] for helpful discussions.
[Funding information and HPC allocation acknowledgements go here.]
This work used resources of [HPC centre].
\end{acknowledgements}

%% ============================================================
%% REFERENCES
%% ============================================================
\bibliographystyle{copernicus}
\bibliography{references}

\end{document}
